%-------------------------------------------------------------------------------
% Sophia Tsang -- Resume
% Compile with pdfLaTeX (e.g., on Overleaf). Single-column, ATS-friendly.
%-------------------------------------------------------------------------------
\documentclass[letterpaper,11pt]{article}

\usepackage{accsupp}
\usepackage[empty]{fullpage}
\usepackage[margin=0.5in]{geometry}
\usepackage{titlesec}
\usepackage[usenames,dvipsnames]{color}
\usepackage{enumitem}
\usepackage[hidelinks]{hyperref}
\usepackage[english]{babel}
\usepackage{tabularx}
\usepackage[T1]{fontenc}
\usepackage[utf8]{inputenc}
% \usepackage{lmodern}
\usepackage[lining]{ebgaramond}

% Make the PDF text machine-readable for applicant tracking systems
\input{glyphtounicode}
\pdfgentounicode=1

\pagestyle{empty}
\raggedright
\setlength{\tabcolsep}{0in}
\urlstyle{same}

% Section formatting
\titleformat{\section}
  {\vspace{-8pt}\scshape\raggedright\largish\bfseries}
  {}{0em}{}[\color{black}\titlerule\vspace{-6pt}]

\usepackage[normalem]{ulem}   % normalem keeps \emph as italics
\setlength{\ULdepth}{2pt}     % distance from baseline down to the underline
\renewcommand{\ULthickness}{0.4pt}  % line thickness

\newcommand{\link}[2]{\href{#1}{\uline{#2}}}

\newcommand{\largish}{\fontsize{14pt}{14pt}\selectfont}

% Line spacing for bullets and education text (second number = spacing)
\newcommand{\bodyspacing}{\fontsize{11pt}{14pt}\selectfont}

% Entry heading: left text, right text, on one line
\newcommand{\entry}[4]{%
  \item
  \begin{tabular*}{\textwidth}[t]{l@{\extracolsep{\fill}}r}
    #1 & #2 \\
    #3 & #4 \\
  \end{tabular*}\vspace{-6pt}}

\newcommand{\education}[4]{\vspace{2pt}%
  \item #1 \hfill #2 \\
  #3 \hfill #4 \\}

\newcommand{\project}[2]{%
  \item
  \begin{tabular*}{\textwidth}{l@{\extracolsep{\fill}}r}
    #1 & #2 \\
  \end{tabular*}\vspace{-6pt}}

\newcommand{\bullets}{\begin{itemize}[leftmargin=0.18in, itemsep=0pt, parsep=0pt, topsep=3pt, partopsep=0pt, label=\textbullet]}
\newcommand{\bulletsend}{\end{itemize}\vspace{-3pt}}
\newcommand{\entries}{\begin{itemize}[leftmargin=0in, label={}, itemsep=0pt]}
\newcommand{\entriesend}{\end{itemize}}

%-------------------------------------------------------------------------------
\begin{document} \fontsize{11pt}{14pt}\selectfont

%----------HEADING----------
\begin{center}
  {\fontsize{19pt}{19pt} \textbf{Sophia Tsang}} \\ \vspace{3pt}
   609-955-2453 $|$
  \href{mailto:xt73@cornell.edu}{\underline{xt73@cornell.edu}} $|$
  \link{https://sophtsang.github.io}{sophtsang.github.io} $|$
  \link{https://github.com/sophtsang}{github.com/sophtsang} $|$
  \link{https://linkedin.com/in/xueqing-tsang}{linkedin.com/in/xueqing-tsang}
\end{center}

%----------EDUCATION----------
\section{Education}
\entries
    \education{\textbf{Cornell University}, College of Engineering}{\textbf{Expected May 2027}}
            {\textit{Bachelor of Science in Computer Science, Minor in Robotics}}{\textbf{GPA: 3.923}}
    \textbf{\textit{Relevant Courses:}} Systems Programming, Formal Verification, Frontiers of Computer Vision, Intro to ML,
        Foundations of Robotics, Robot Learning, Robot Perception, Intro to Game Architecture,
        Analysis of Algorithms, Differential Equations
\entriesend

%----------EXPERIENCE----------
\section{Experience}
\entries
  \entry{\large \scshape \textbf{MIT Lincoln Laboratory}}{\textbf{Lexington, MA}}
        {\textbf{Software Engineering and Robotics Intern}  $|$ Python, PyTorch, TensorFlow, Docker, ROS2, VIO}{\textbf{May 2026 -- Aug 2026}}
  \bullets
    \item Engineered an end-to-end 3D terrain reconstruction pipeline for UAV imagery, 
      adapting \link{https://arxiv.org/html/2412.13047v1}{\textit{Earth Observation Gaussian Splatting}}
      and shadow mapping to generate depth maps, achieving
      \textbf{92.0\% geometric accuracy} against LiDAR ground-truth terrain
    \item Integrated camera-derived depth and terrain maps into the3D scene-graph codebase
      and evaluation framework, enabling an alternative, physically grounded depth-estimation method 
      that works with limited aerial perspectives for terrain reconstruction
    \item Architected asynchronous ROS2 processing with multiprocessing to decouple camera-stream ingestion from background depth estimation, 
      enabling \textit{non-blocking data processing and continuous 
      depth-map generation}
  \bulletsend

  \entry{\large \scshape \link{https://www.cornellelectricvehicles.org/}{\textbf{Cornell Electric Vehicles}}}{\textbf{Ithaca, NY}}
        {\textbf{Perception Lead, Autonomy Engineer}  $|$ Rust, Python, C++, PyTorch, ROS2, Docker, Cargo}{\textbf{Oct 2024 -- Present}}
  \bullets
    \item Lead a 5-person Perception subsystem responsible for environmental sensing and world representation, overseeing feature development, 
      project ownership, engineering milestones, and cross-team integration for vehicle perception systems
    \item Engineered \textit{C++/Rust LiDAR driver}, owning the full data-processsing pipeline from UDP packet ingestion and
      IMU-based ego-motion compensation to \textit{world-aligned 3D point cloud generation}, \textbf{reducing static-obstacle drift} 
      from \textbf{0.92m to 0.04m}
    \item Optimized \link{https://www.mdpi.com/2076-3417/10/23/8534}{\textit{two-layer-graph clustering}}
      on \textbf{80k+ points} using azimuth-binning and range-adaptive thresholds, reducing scan latency to
      \textbf{20--30ms} and far-range over-merging from \textbf{99\% to <10\%} for \textit{real-time obstacle detection and mapping}
    \item Developed \textit{occlusion-tolerant multi-object tracking} and point-cloud-to-\textit{occupancy-grid} generation, maintaing 
      persistent obstacle identities and environmental state for \textit{downstream planning and decision-making}
    % \item Implemented a \textit{real-time} point cloud to occupancy grid algorithm that constructs accurate
    %   representations of known free-space and obstacle regions around the vehicle,
    %   \textit{enabling accurate and continuous spatial mapping} for path planning
    % \item Integrated camera–LiDAR fusion to combine RGB imagery with precise depth measurements for downstream algorithms
    % \item Simulated LiDAR scans with a physics-based measurement model for \textit{testing perception algorithms} in variable environments
  \bulletsend

  \entry{\large \scshape \textbf{Noah Snavely Research Group}}{\textbf{Roosevelt Island, NY}}
        {{\textbf{Research Intern}} $|$ PyTorch, COLMAP, CUDA, Gemini, GCP, Flask, Render, React, Three.js, WebGL}{\textbf{June 2025 -- Jan 2026}}
  \bullets
    \item Engineered an automated geospatial data-collection pipeline integrating Google Maps API
      to identify landmark boundaries, interpolate candidate street-level viewpoints, and assemble georeferences 
      imagery datasets across \textbf{4000+ landmarks}
    \item Developed a \textit{geo-referencing pipeline} to recover absolute metric scale from SfM reconstructions using 
      geographically tagged imagery and optimized viewpoint selection, co-authoring the ECCV 2026 paper,
      \link{https://arxiv.org/pdf/2606.02379}{``Honey, I Shrunk the Arc de Triomphe''}
    % \item Expanded \link{https://doppelgangers25.github.io/doppelgangers_plusplus/}{Doppelgangers++}
    %   into a vision-language network for \textit{end-to-end 3D reconstruction}, enabling visual 
    %   disambiguation between distinct but visually similar surfaces as a more robust alternative to classical SfM
    % \item Built a web viewer that runs on React, WebGL, and Flask to render
    %   \textbf{4k+ doppelganger-pruned 3D reconstructions}
  \bulletsend

\entriesend

%----------PROJECTS----------
\section{Projects}
\entries

  \project{\link{https://sophtsang.github.io/}{\textbf{Kernel::BSOD: OS Simulator}}  $|$ {Rust, WebAssembly, JavaScript, React, Three.js, Cargo}}{\textbf{Aug 2026 -- Present}}
  \bullets
    \item Built a Rust process-management subsystem compiled to WebAssembly, implementing process control blocks,
      process lifecycle state machines, and zombie-process reaping to manage application instances through OS-inspired 
      process states
    \item Implemented a binary space partitioning (BSP) window manager in Rust, integrating process lifecycle events with UI rendering to dynamically allocate window geometry and reposition 
      tiled windows based on user-triggered focus
  \bulletsend

  \project{\link{https://gdiac.cs.cornell.edu/gdiac/showcase/games/re_color/}{\textbf{Re:Color}} $|$ {Java, LibGDX, Agile Project Management}}{\textbf{Jan 2026 -- May 2026}}
  \bullets
    \item Designed a color-based gameplay system with 7 other developers, enabling players to
      recolor the world by painting and absorbing color from the 2D platformer environment, and
      shipped the game \textit{live} at GDIAC Showcase to \textbf{400+ attendees}
    \item Developed and shipped core player features and systems, including behavior-tree-driven 
      player state management, animation synchronization, movement physics refinement, and 
      a custom physics-based paint splatter rendering engine
  \bulletsend

\entriesend

%----------SKILLS----------
\section{Skills}
\begin{itemize}[leftmargin=0in, label={}]
  {\item{
    \textbf{\textit{Languages:}} Python, Rust, C/C++, C\#, Java, Rocq, OCaml, JavaScript/TypeScript, SQL,
      HTML, CSS, Verilog \\
    \textbf{\textit{Frameworks/Libraries:}} PyTorch, NumPy, React.js, Three.js, WebAssembly, Flask,
      OpenCV, sqlite3, WebGL, Matplotlib \\
    \textbf{\textit{Tools:}} Git, CI/CD, Cargo, ROS2, Linux, Docker, VSCode, LibGDX, Unity, Render, GCP,
      Figma, Fusion 360, Slurm
  }}
\end{itemize}

\end{document}